\chapter{绪论}

\section{研究背景}

跨媒介职业转型已经成为东亚娱乐产业中的常见现象。传统纸媒时代，职业晋升更多依赖杂志封面、广告代言与电视剧曝光；进入视频平台时代后，艺人需要在更高频的互动环境中持续解释自我、修复争议并管理公众期待。因此，职业转型不再只是“工作类型改变”，而是平台逻辑、内容机制与公众评价共同变化的结果。

\section{案例价值}

佐佐木希的公开职业轨迹兼具跨度长、阶段清晰和资料丰富三项特点。她从模特出道，逐步进入影视和综艺，再到视频平台中的自我表达，能够为研究“同一主体多次转型”提供足够清晰的结构。相关事件与舆论转折也具有公开来源，适合在不暴露保密材料的前提下构建完整的硕士论文示例。

\section{研究目的与内容}

本示例的任务有两项。其一，基于公开资料概括跨媒介职业转型中的平台选择和形象韧性机制。其二，以可分享正文验证 `thesis-ahnu-master` 的封面、题名页、中英文摘要、目录、缩略词表、章节正文和后置材料等排版链路是否稳定。
